\section{Reproducibility statement}
\label{sec:reproducibility}

Written early on purpose: it disciplines the rest.

\textbf{Which figures are recomputable, and from what.} Activation rates,
outcome counts and all dispersion figures come from
\texttt{studies/run\_ledger.csv}. Retrieval figures come from the labelled
query sets in \texttt{tests/} and are reproduced by
\texttt{internal/tools/source\_index\_baseline.py}. Prompt provenance is
regenerated by \texttt{make promptmap} from the live graph objects, and a test
fails if a prompt moves out from under its citation.

\textbf{The line historical data falls on.} Figures computed from runs
predating the fixes in \cref{sec:threats} inherit those defects; the ledger's
success label is unreliable before them. The ablation runs on the fixed system,
and the two sets of figures are never mixed in one table.

\textbf{Known reproducibility defects.} The anchor study that every variance
figure is drawn from is not committed to the repository. No seed or temperature
control exists. The Abaqus corpus cannot be redistributed, so that tool's
quality is unmeasured by anyone without a licensed installation.

\textbf{What is pinned.} The model is named per arm. The resolved configuration
each run actually used is recorded per run rather than inferred from the
study's configuration file.
